\documentclass[11pt]{article}
\usepackage[a4paper,margin=22mm]{geometry}
\usepackage[no-math]{fontspec}
\setmainfont{DejaVu Serif}
\usepackage{amsmath,amssymb,amsthm,mathtools,graphicx,xcolor,xurl,xspace,enumerate}
\usepackage[shortlabels]{enumitem}
\setlength{\emergencystretch}{4em}
\providecommand{\smaller}{\small}
\providecommand{\loccit}{\textup{loc.\ cit.}}
\providecommand{\locscit}{\textup{locc.\ cit.}}
\providecommand{\ibid}{\textup{ibid.}}
\providecommand{\cf}{\textup{cf.}}
\providecommand{\resp}{\textup{resp.}}
\providecommand{\eg}{\textup{e.g.}}
\providecommand{\ie}{\textup{i.e.}}
\providecommand{\vs}{\textup{vs.}}
%
% retained author preamble lines below (source lines 1-119, publisher class and
% publisher meta commands removed); tex of the card body is delivered verbatim.
%

\usepackage[scr=rsfs,cal=euler]{mathalfa}
\newcommand{\llbracket}{[\![}
\newcommand{\rrbracket}{]\!]}
\newcommand{\sbullet}{{\scriptscriptstyle\bullet}}
\newcommand{\csbullet}{{\raisebox{1pt}{$\sbullet$}}}
\newcommand\mto{\mathchoice{\longmapsto}{\mapsto}{\mapsto}{\mapsto}}
\newcommand\hto{\mathrel{\lhook\joinrel\to}}
\newcommand\To[1]{\mathchoice{\xrightarrow{\kern1pt#1\kern1pt}}{\stackrel{#1}{\longrightarrow}}{}{}}
\newcommand{\quand}{\quad\text{and}\quad}
\let\epsilon\varepsilon

\usepackage{tikz-cd}
\usepackage[all,cmtip]{xy}
\newdir{ >}{{}*!/-5pt/\dir{>}}

\theoremstyle{plain}
\newtheorem{theorem}{Theorem}[section]
\newtheorem{proposition-definition}{Proposition/Definition}[section]
\newtheorem{theorem-definition}{Theorem/Definition}[section]
\newtheorem*{maintheorem}{Main Theorem}
\newtheorem*{proposition-definition-intro}{Proposition/Definition}
\newtheorem*{definition-intro}{Definition}
\newtheorem*{cor-intro}{Corollary}
\newtheorem*{prop-intro}{Proposition}
\newtheorem{proposition}[theorem]{Proposition}		
\newtheorem{corollary}[theorem]{Corollary}
\newtheorem{lemma}[theorem]{Lemma}
\newtheorem*{conjecture}{Conjecture}
\newtheorem*{conjecture-intro}{Conjecture}
\newtheorem{question}[theorem]{Question}
\newtheorem{claim}[theorem]{Claim}
\newtheorem*{theoremMeta}{Main theorem}
\newtheorem{thm}{Theorem}
\renewcommand*{\thethm}{\Alph{thm}}
\newtheorem{prop}[thm]{Proposition}

\theoremstyle{definition}
\newtheorem{definition}[theorem]{Definition}
\newtheorem{construction-definition}[theorem]{Construction/Definition}
\newtheorem{notation}[theorem]{Notation}
\newtheorem*{problem2prime}{Problem 2'}
\newtheorem{defthm}[thm]{Definition / Theorem}
\newtheorem{problem-Deligne}{Problem}
\newtheorem{remark}[theorem]{Remark}
\newtheorem*{remark-intro}{Remark}
\newtheorem{example}[theorem]{Example}
\newtheorem*{construction}{Construction}


\newcommand{\Abold}{\mathbf{A}}
\newcommand{\Bbold}{\mathbf{B}}
\newcommand{\Cbold}{\mathbf{C}}
\newcommand{\Dbold}{\mathbf{D}}
\newcommand{\Ebold}{\mathbf{E}}
\newcommand{\Kbold}{\mathbf{K}}
\newcommand{\Pbold}{\mathbf{P}}

\newcommand{\NBbb}{\mathbb N}
\newcommand{\PBbb}{\mathbb P}
\newcommand{\QBbb}{\mathbb Q}
\newcommand{\ZBbb}{\mathbb Z}

\newcommand{\Acal}{\mathcal A}
\newcommand{\Bcal}{\mathcal B}
\newcommand{\Ccal}{\mathcal C}
\newcommand{\Dcal}{\mathcal D}
\newcommand{\Ecal}{\mathcal E}
\newcommand{\Gcal}{\mathcal G}
\newcommand{\Ical}{\mathcal I}
\newcommand{\Ocal}{\mathcal O}
\newcommand{\Pcal}{\mathcal P}
\newcommand{\Rcal}{\mathcal R}
\newcommand{\Scal}{\mathcal S}
\newcommand{\Ucal}{\mathcal U}
\newcommand{\Vcal}{\mathcal V}
\newcommand{\Wcal}{\mathcal W}

\newcommand{\Cfrak}{\mathfrak C}
\newcommand{\Hfrak}{\mathfrak H}

\newcommand{\cfrak}{\mathfrak c}
\newcommand{\chfrak}{\mathfrak {ch}}
\newcommand{\sfrak}{\mathfrak s}
\newcommand{\tdfrak}{\mathfrak {td}}

\newcommand{\bis}{\prime\prime}
\newcommand{\tris}{\prime\prime\prime}
\newcommand{\fppf}{\emph{fppf}\xspace}

\newcommand{\Pic}{\mathrm{Pic}}
\newcommand{\prim}{\mathrm{prim}}

\DeclareMathOperator{\Aut}{Aut}
\DeclareMathOperator{\ch}{ch}
\DeclareMathOperator{\CH}{CH}
\DeclareMathOperator{\CHfrak}{\mathfrak{CH}}
\DeclareMathOperator{\Gr}{Gr}
\DeclareMathOperator{\Hom}{Hom}
\DeclareMathOperator{\id}{id}
\DeclareMathOperator{\Imag}{Im}
\DeclareMathOperator{\Picfr}{\mathfrak{Pic}}
\DeclareMathOperator{\iso}{iso}
\DeclareMathOperator{\qiso}{quasi-iso}
\DeclareMathOperator{\rk}{rk}
\DeclareMathOperator{\Spec}{Spec}
\DeclareMathOperator{\Sym}{Sym}
\DeclareMathOperator{\td}{td}
\DeclareMathOperator{\Vect}{Vect}

\let\ov\overline

\usepackage{hyperref}
\usepackage{cleveref}
\crefname{originalproposition}{original proposition}{original propositions}
\Crefname{originalproposition}{Original proposition}{Original propositions}
\makeatletter
\newlabel{prop:generalpropertiesdeligneproducts}{{6.1}{source}{}{sourcepointer-1}{}}
\newlabel{prop:generalpropertiesdeligneproducts@cref}{{[originalproposition][0][]6.1}{[1][0][]source}{}{}{}}
\makeatother
\crefname{originalproposition}{original proposition}{original propositions}
\Crefname{originalproposition}{Original proposition}{Original propositions}
\makeatletter
\newlabel{item:deligne-pairing-restriction}{{deligne-pairing-restriction}{source}{}{sourcepointer-2}{}}
\newlabel{item:deligne-pairing-restriction@cref}{{[originalproposition][0][]deligne-pairing-restriction}{[1][0][]source}{}{}{}}
\makeatother
\crefname{originalcorollary}{original corollary}{original corollarys}
\Crefname{originalcorollary}{Original corollary}{Original corollarys}
\makeatletter
\newlabel{cor:compprojformula}{{6.13}{source}{}{sourcepointer-3}{}}
\newlabel{cor:compprojformula@cref}{{[originalcorollary][0][]6.13}{[1][0][]source}{}{}{}}
\makeatother
\crefname{originalproposition}{original proposition}{original propositions}
\Crefname{originalproposition}{Original proposition}{Original propositions}
\makeatletter
\newlabel{prop:projformuladeli}{{6.11}{source}{}{sourcepointer-4}{}}
\newlabel{prop:projformuladeli@cref}{{[originalproposition][0][]6.11}{[1][0][]source}{}{}{}}
\makeatother
\crefname{originalproposition}{original proposition}{original propositions}
\Crefname{originalproposition}{Original proposition}{Original propositions}
\makeatletter
\newlabel{prop:split}{{5.39}{source}{}{sourcepointer-5}{}}
\newlabel{prop:split@cref}{{[originalproposition][0][]5.39}{[1][0][]source}{}{}{}}
\makeatother
\crefname{originalproposition}{original proposition}{original propositions}
\Crefname{originalproposition}{Original proposition}{Original propositions}
\makeatletter
\newlabel{Prop:2ndsplitting}{{5.41}{source}{}{sourcepointer-6}{}}
\newlabel{Prop:2ndsplitting@cref}{{[originalproposition][0][]5.41}{[1][0][]source}{}{}{}}
\makeatother
\crefname{originalproposition}{original proposition}{original propositions}
\Crefname{originalproposition}{Original proposition}{Original propositions}
\makeatletter
\newlabel{prop:elvira}{{6.14}{source}{}{sourcepointer-7}{}}
\newlabel{prop:elvira@cref}{{[originalproposition][0][]6.14}{[1][0][]source}{}{}{}}
\makeatother
\crefname{originallemma}{original lemma}{original lemmas}
\Crefname{originallemma}{Original lemma}{Original lemmas}
\makeatletter
\newlabel{lemma:elementary-properties-segre}{{7.2}{source}{}{sourcepointer-8}{}}
\newlabel{lemma:elementary-properties-segre@cref}{{[originallemma][0][]7.2}{[1][0][]source}{}{}{}}
\makeatother
\crefname{originalproposition}{original proposition}{original propositions}
\Crefname{originalproposition}{Original proposition}{Original propositions}
\makeatletter
\newlabel{item:elpropsef-0}{{elpropsef-0}{source}{}{sourcepointer-9}{}}
\newlabel{item:elpropsef-0@cref}{{[originalproposition][0][]elpropsef-0}{[1][0][]source}{}{}{}}
\makeatother
\crefname{originalproposition}{original proposition}{original propositions}
\Crefname{originalproposition}{Original proposition}{Original propositions}
\makeatletter
\newlabel{item:elpropseg-2}{{7.2(2)}{source}{}{sourcepointer-10}{}}
\newlabel{item:elpropseg-2@cref}{{[originalproposition][0][]7.2(2)}{[1][0][]source}{}{}{}}
\makeatother
\crefname{originalproposition}{original proposition}{original propositions}
\Crefname{originalproposition}{Original proposition}{Original propositions}
\makeatletter
\newlabel{item:elpropseg-3}{{7.2(3)}{source}{}{sourcepointer-11}{}}
\newlabel{item:elpropseg-3@cref}{{[originalproposition][0][]7.2(3)}{[1][0][]source}{}{}{}}
\makeatother
\crefname{originalremark}{original remark}{original remarks}
\Crefname{originalremark}{Original remark}{Original remarks}
\makeatletter
\newlabel{rmk:taut-proj-form}{{7.4}{source}{}{sourcepointer-12}{}}
\newlabel{rmk:taut-proj-form@cref}{{[originalremark][0][]7.4}{[1][0][]source}{}{}{}}
\makeatother
\crefname{originalremark}{original remark}{original remarks}
\Crefname{originalremark}{Original remark}{Original remarks}
\makeatletter
\newlabel{rmk:clarification-restriction}{{6.2}{source}{}{sourcepointer-13}{}}
\newlabel{rmk:clarification-restriction@cref}{{[originalremark][0][]6.2}{[1][0][]source}{}{}{}}
\makeatother
\crefname{originalequation}{original equation}{original equations}
\Crefname{originalequation}{Original equation}{Original equations}
\makeatletter
\newlabel{eq:definition-chern-classes}{{definition-chern-classes}{source}{}{sourcepointer-14}{}}
\newlabel{eq:definition-chern-classes@cref}{{[originalequation][0][]definition-chern-classes}{[1][0][]source}{}{}{}}
\makeatother
\title{Functorial properties of the intersection bundles associated with Chern power series}
\author{Dennis Eriksson; Gerard Freixas i Montplet}
\date{doi:10.5802/jep.254; CC BY 4.0}
\begin{document}
\maketitle
\noindent\textit{Editorial note.} All delivered mathematical text is the author's own, copied from the pinned publisher source of this unit; no statement or proof step was rewritten, reordered or removed. The counted statement is the labelled proposition environment at source lines 2450-2492 (with its six numbered items) and the counted proof is the proof environment at source lines 2494-2555.\par\medskip
\noindent\textit{Editorial note.} The preamble of the delivered fragment is generated for this unit: the publisher class is replaced by article plus the ams packages and the author's own package, macro and theorem-declaration lines are retained in source order; publisher meta commands are removed.\par\medskip
\noindent\textit{Editorial note.} The results, items, remarks and labelled displays invoked by the counted text are declared as non-load-bearing external pointers rather than restated, except for one labelled equation of the same article delivered as context; the citation of the source text to a book outside the corpus is resolved through the publisher reference list retained for this unit.\par\medskip
\noindent\textit{Allowed standard prerequisites (editorial).}\begin{itemize}
\item Basic properties of Deligne pairings and projection formulas. The author's proposition on the basic properties of Deligne pairings (printed number 6.1) and the proposition supplying the projection formula (printed number 6.11). Give the projection formula isomorphisms of the first assertion and the compatibility rules for their composition.
\item Splitting principles. The author's two splitting propositions (printed numbers 5.39 and 5.41). Reduce the constructions to the case of line bundles, where the verification is explicit.
\item Elementary properties and the restriction isomorphism. The author's lemma on the elementary properties (printed number 7.2), the earlier proposition on Deligne pairings (printed number 6.1) and the author's remark on the tautological projection formula (printed number 7.4). Supply the multilinearity, the determinant rule and the restriction isomorphism used in the first Chern class, rank triviality and restriction constructions.
\item Birational invariance. The author's proposition on birational invariance (printed number 6.14). Yields the last construction once the monomial is written in terms of Segre classes of vector bundles.
\end{itemize}\medskip
\ifdefined\hypertarget\hypertarget{sourcepointer-1}{}\fi
\noindent\textit{Source reference.} Original proposition 6.1: sources/jep-2024-254-deligne-riemann-roch-and-intersection-bundles/pinned-original-source.tex, source lines 1965--1986. Prerequisite pointer invoked by the counted proof for the basic properties of Deligne pairings, in particular the restriction isomorphism employed in the projections.\par\medskip
\ifdefined\hypertarget\hypertarget{sourcepointer-2}{}\fi
\noindent\textit{Source reference.} Original proposition deligne-pairing-restriction: sources/jep-2024-254-deligne-riemann-roch-and-intersection-bundles/pinned-original-source.tex, source lines 1975--1975. Prerequisite pointer invoked by the counted proof: the item recording the behaviour of a Deligne pairing under restriction.\par\medskip
\ifdefined\hypertarget\hypertarget{sourcepointer-3}{}\fi
\noindent\textit{Source reference.} Original corollary 6.13: sources/jep-2024-254-deligne-riemann-roch-and-intersection-bundles/pinned-original-source.tex, source lines 2233--2241. Prerequisite pointer invoked by the counted statement and proof: the composition rule for projection formulas that the constructed isomorphisms must satisfy.\par\medskip
\ifdefined\hypertarget\hypertarget{sourcepointer-4}{}\fi
\noindent\textit{Source reference.} Original proposition 6.11: sources/jep-2024-254-deligne-riemann-roch-and-intersection-bundles/pinned-original-source.tex, source lines 2169--2187. Prerequisite pointer invoked by the counted proof: it supplies the projection formula from which the first construction formally follows.\par\medskip
\ifdefined\hypertarget\hypertarget{sourcepointer-5}{}\fi
\noindent\textit{Source reference.} Original proposition 5.39: sources/jep-2024-254-deligne-riemann-roch-and-intersection-bundles/pinned-original-source.tex, source lines 1879--1891. Prerequisite pointer invoked by the counted proof as the first splitting principle used in the reduction to line bundles.\par\medskip
\ifdefined\hypertarget\hypertarget{sourcepointer-6}{}\fi
\noindent\textit{Source reference.} Original proposition 5.41: sources/jep-2024-254-deligne-riemann-roch-and-intersection-bundles/pinned-original-source.tex, source lines 1899--1900. Prerequisite pointer invoked by the counted proof as the second splitting principle.\par\medskip
\ifdefined\hypertarget\hypertarget{sourcepointer-7}{}\fi
\noindent\textit{Source reference.} Original proposition 6.14: sources/jep-2024-254-deligne-riemann-roch-and-intersection-bundles/pinned-original-source.tex, source lines 2246--2263. Prerequisite pointer invoked by the counted proof in the treatment of birational invariance.\par\medskip
\ifdefined\hypertarget\hypertarget{sourcepointer-8}{}\fi
\noindent\textit{Source reference.} Original lemma 7.2: sources/jep-2024-254-deligne-riemann-roch-and-intersection-bundles/pinned-original-source.tex, source lines 2393--2414. Prerequisite pointer invoked by the counted proof for the elementary properties of the pairings used in the first Chern class and rank triviality constructions.\par\medskip
\ifdefined\hypertarget\hypertarget{sourcepointer-9}{}\fi
\noindent\textit{Source reference.} Original proposition elpropsef-0: sources/jep-2024-254-deligne-riemann-roch-and-intersection-bundles/pinned-original-source.tex, source lines 2396--2396. Prerequisite pointer invoked by the counted proof: the item of the elementary properties recording the isomorphism of pairings induced by an isomorphism of bundles.\par\medskip
\ifdefined\hypertarget\hypertarget{sourcepointer-10}{}\fi
\noindent\textit{Source reference.} Original proposition 7.2(2): sources/jep-2024-254-deligne-riemann-roch-and-intersection-bundles/pinned-original-source.tex, source lines 2406--2406. Prerequisite pointer invoked by the counted proof: the multilinearity item of the elementary properties.\par\medskip
\ifdefined\hypertarget\hypertarget{sourcepointer-11}{}\fi
\noindent\textit{Source reference.} Original proposition 7.2(3): sources/jep-2024-254-deligne-riemann-roch-and-intersection-bundles/pinned-original-source.tex, source lines 2407--2407. Prerequisite pointer invoked by the counted proof: the item of the elementary properties on the determinant of an extension.\par\medskip
\ifdefined\hypertarget\hypertarget{sourcepointer-12}{}\fi
\noindent\textit{Source reference.} Original remark 7.4: sources/jep-2024-254-deligne-riemann-roch-and-intersection-bundles/pinned-original-source.tex, source lines 2440--2446. Prerequisite pointer invoked by the counted proof for the identification of the last term of the Whitney computation.\par\medskip
\ifdefined\hypertarget\hypertarget{sourcepointer-13}{}\fi
\noindent\textit{Source reference.} Original remark 6.2: sources/jep-2024-254-deligne-riemann-roch-and-intersection-bundles/pinned-original-source.tex, source lines 1992--2002. Prerequisite pointer invoked by the counted proof for the convention on restriction isomorphisms.\par\medskip
\ifdefined\hypertarget\hypertarget{sourcepointer-14}{}\fi
\noindent\textit{Source reference.} Original equation definition-chern-classes: sources/jep-2024-254-deligne-riemann-roch-and-intersection-bundles/pinned-original-source.tex, source lines 2366--2366. Prerequisite pointer invoked by the counted proof: the defining recurrence of the Chern classes of a line bundle.\par\medskip
\begin{proposition} \label{prop:properties-chern-bundles}
Let $f\colon X\to S$ be a morphism satisfying the condition $(C_{n})$, and let $P$, $P'$ denote positive Chern power series. The following constructions are functorial:
\begin{enumerate}

\item\label{item:projection-fromula} (Projection formulas) Let $g\colon X^{\prime}\to S$ and $h \colon X' \to X$ satisfy the conditions $(C_{n+n'})$ and $(C_{n'})$, respectively, and $P$, $P'$ have pure degrees with $\deg P + \deg P'$ being $n'+n+1$. Then there are canonical projection formula isomorphisms
\begin{displaymath}
\langle h^\ast P \cdot P' \rangle_{X'/S} \simeq
\begin{cases}
\langle P\cdot \langle P' \rangle_{X'/X} \rangle_{X/S}, & \text{if }\; \deg P^{\prime}=n^{\prime}+1.\\
\langle P \rangle_{X/S}^{\kappa}, & \text{if }\; \deg P^{\prime} = n^{\prime}. \\
\Ocal_S, & \text{if }\; \deg P^{\prime} < n^{\prime}.
\end{cases}
\end{displaymath}
In the first case, we wrote $\langle P' \rangle_{X'/X}$ instead of $\cfrak_{1}(\langle P' \rangle_{X'/X})$ for simplicity. In the second case, $\kappa$ denotes the relative degree $\int_{X'/X} P'$, and is then assumed to be constant. In the third case, we interpret $\Ocal_S$ as a constant functor. These isomorphisms satisfy the analogs of Corollary \ref{cor:compprojformula}.

\item \label{whitneyiso} (Whitney isomorphism) Let $0 \to E^{\prime}\to E \to E^{\bis} \to 0$ be a short exact sequence of vector bundles on $X$. Then there is a canonical isomorphism
\begin{displaymath}
\langle P \cdot \cfrak_{k}(E)\cdot P'\rangle_{X/S} \to \bigotimes_{i=0}^{k} \langle P \cdot \cfrak_{i}(E^{\prime}) \cdot \cfrak_{k-i}(E^{\bis})\cdot P' \rangle_{X/S},
\end{displaymath}
in a way that is compatible with admissible filtrations.

\item\label{item:first-chern-determinant} (First Chern class isomorphism) Let $E$ be a vector bundle on $X$. Then, there is a canonical isomorphism
\begin{displaymath}
\langle P\cdot \cfrak_{1}(E)\cdot P'\rangle_{X/S} \simeq \langle P\cdot \cfrak_{1}(\det E)\cdot P'\rangle_{X/S},
\end{displaymath}
that is compatible with the Whitney isomorphism.
\item \label{item:Chern-rank}(Rank triviality) Let $E$ be a vector bundle on $X$ and $q$ an integer such that $q > \rk E$. Then the intersection bundle
\begin{displaymath}
\langle P \cdot \cfrak_q(E) \cdot P' \rangle_{X/S}
\end{displaymath}
is canonically isomorphic to the constant functor $\Ocal_S$. Moreover, the isomorphism is compatible with the Whitney isomorphism.
\item\label{item:restriction-isomorphism} (Restriction isomorphism) Let $\sigma$ be a regular section of a vector bundle $E$, of~constant rank $r$, whose zero locus $Y$, possibly empty, is flat over $S$. Then there is a canonical isomorphism
\begin{displaymath}
\langle P \cdot \cfrak_{r}(E) \cdot P' \rangle_{X/S} \to \langle P|_Y \cdot P'|_Y \rangle_{Y/S},
\end{displaymath}
where we adopt the same convention as in Proposition \ref{prop:generalpropertiesdeligneproducts}\eqref{item:deligne-pairing-restriction} (\cf Remark \ref{rmk:clarification-restriction}).
\item \label{item:birational-invariance} (Birational invariance) Let $g: \widetilde{X} \to S$ be a morphism satisfying condition $(C_n)$. Assume that there exists a morphism $\pi\colon \widetilde{X} \to X$ and a quasi-compact open immersion $U \to X$, such that $\pi^{-1}(U) \to U$ is an isomorphism and $U$ is fiberwise dense in $X$.
Then there is a canonical isomorphism
\begin{displaymath}
\langle P \rangle_{X/S}\simeq \langle \pi^\ast P \rangle_{\widetilde{X}/S}.
\end{displaymath}
\end{enumerate}
\end{proposition}

\begin{proof}

The projection formula formally follows from the construction of Segre and Chern classes together with Proposition \ref{prop:projformuladeli}, as in
\cite[\textsection V.4.2]{Elkikfib}. Notice that in \loccit, $X$ is only assumed to be connected in the general case, for simplicity. The analog of Corollary \ref{cor:compprojformula} similarly follows.

The Whitney isomorphism is constructed similarly as in \cite[\textsection V.4.8]{Elkikfib}. Before recalling it, we remark it is only defined therein whenever the vector bundles have positive ranks. In the case $\alpha: E^{\prime} \to E$ is an isomorphism, so that $E^{\bis}=0$, we define the Whitney isomorphism as $[\alpha]^{-1}$, where $[\alpha]$ is the isomorphism of intersection bundles induced by $\alpha$ as in Lemma \ref{lemma:elementary-properties-segre}\eqref{item:elpropsef-0}. Likewise, if $\alpha: E\to E^{\bis}$ is an isomorphism and $E^{\prime}=0$, we define the Whitney isomorphism as $[\alpha]$.

The construction in the general case reduces to a Whitney-type isomorphism for Segre classes. By the splitting principle in Proposition \ref{prop:split} and Proposition \ref{Prop:2ndsplitting}, it~is enough to describe this in the case when $E^{\prime}=L$ is a line bundle. Let $r$ be the rank of~$E$. We may suppose that $r\geq 2$. Denote by $\pi\colon\PBbb(E)\to X$ the projection map, and~$\Ocal(1)$ the tautological bundle on $\PBbb(E)$. We can write, successively replacing a single~$\Ocal(1)$ by $\Ocal(1) \otimes \pi^{\ast} L^{\vee} \otimes \pi^\ast L $ and iterating the process $k$ times :
\begin{equation}\label{eq:SegreWhitney}
\begin{split}
\langle P \cdot \sfrak_k(E) &{}\cdot P' \rangle = \langle \pi^{\ast}P\cdot \cfrak_1(\Ocal(1))^{r+k-1} \cdot \pi^{\ast}P' \rangle \\
&= \langle \pi^{\ast}P\cdot \cfrak_1(\Ocal(1) \otimes \pi^{\ast}L^\vee) \cdot \cfrak_1(\Ocal(1))^{r+k-2}\cdot \pi^{\ast}P' \rangle\\
&\hspace*{2cm}\otimes \langle \pi^{\ast}P\cdot \cfrak_1(\pi^{\ast}L)\cdot\cfrak_1(\Ocal(1))^{r+k-2} \cdot \pi^{\ast}P' \rangle \\
&=\cdots\\
&= \bigotimes_{i=0}^{k-1} \langle \pi^{\ast}P \cdot \cfrak_1(\Ocal(1) \otimes \pi^{\ast} L^\vee) \cdot \cfrak_1(\Ocal(1))^{r+k-2-i} \cdot \cfrak_1(\pi^\ast L)^{i} \cdot \pi^{\ast}P'\rangle \\
&\hspace*{2cm}\otimes \langle \pi^{\ast}P \cdot \cfrak_1(\pi^\ast L)^{k} \cdot \cfrak_1(\Ocal(1))^{r-1} \cdot \pi^{\ast}P'\rangle.
\end{split}
\end{equation}
By Remark \ref{rmk:taut-proj-form}, the last term is identified with $\langle P \cdot \sfrak_k(L) \cdot \sfrak_0(E^{\bis}) \cdot P \rangle$. Also, the bundle $\Ocal(1) \otimes \pi^\ast L^\vee$ admits a natural section $\sigma$, obtained from combining $L \to E$ and $\pi^\ast E \to \Ocal(1)$. Its zero scheme is the relative Cartier divisor $\PBbb(E^{\bis})$ of $\PBbb(E)$. Restricting the bundles along this section and applying Proposition \ref{prop:generalpropertiesdeligneproducts}\eqref{item:deligne-pairing-restriction}, provides the sought Whitney-type isomorphism, again by Remark \ref{rmk:taut-proj-form}. The proof of the compatibility with admissible filtrations is formal, as in \cite[\textsection V.2.3]{Elkikfib}.

The first Chern class isomorphism can be constructed by induction on the rank and the splitting principle. For rank 0, we impose that the isomorphism is the one respecting the trivializations of both sides, see Lemma \ref{lemma:elementary-properties-segre}\eqref{item:elpropseg-2}. The isomorphism in general is then the one that makes the following diagram commute, associated with an exact sequence $0 \to L \to E \to E^{\bis} \to 0$ with $L$ a line bundle:
\begin{equation}\label{eq:firstChernclassisodef}
\begin{array}{c}
\xymatrix@C=.5cm{\langle P\cdot \cfrak_1(E)\cdot P'\rangle \ar[r] \ar[d] & \langle P\cdot\cfrak_1(\det E)\cdot P'\rangle \ar[d] \\
\langle P \cdot \cfrak_1(L) \cdot P'\rangle \otimes \langle P \cdot \cfrak_1(E^{\bis}) \cdot P' \rangle \ar[r] & \langle P \cdot \cfrak_1(L ) \cdot P' \rangle \otimes \langle P \cdot \cfrak_1(\det E^{\bis} ) \cdot P' \rangle.
}
\end{array}
\end{equation}
Here the left vertical isomorphism is the Whitney isomorphism together with Lem\-ma~\ref{lemma:elementary-properties-segre}\eqref{item:elpropseg-3}, the right vertical one follows from $\det E\simeq L\otimes\det E^{\bis}$ together with the multilinearity for $\cfrak_{1}$ (\cf Lemma \ref{lemma:elementary-properties-segre}\eqref{item:elpropseg-2}), and the lower horizontal one is constructed by induction. The compatibility with the Whitney isomorphism is automatic, by the splitting principle and the construction.

The rank triviality is addressed as follows. By the construction of intersection bundles, we can suppose that $E$ has constant non-zero rank. We begin with the case when $E$ is a line bundle $L$. Observe that in terms of intersection bundles since $\Ocal(1)$ identifies with $L$ under the isomorphism $\PBbb(L)\simeq X$, we have
\begin{displaymath}
\langle P \cdot \sfrak_q(L) \cdot P' \rangle = \langle P \cdot \sfrak_1(L)^q \cdot P' \rangle.
\end{displaymath}
From the defining recurrence \eqref{eq:definition-chern-classes}, we deduce that the rank triviality isomorphism is tautological for $q=2$. In general, proceeding by induction on $q$, we find
\begin{equation}\label{eq:Segre-manipulation}
\langle P \cdot \cfrak_{q}(L) \cdot P' \rangle \simeq \langle P \cdot \sfrak_{q}(L) \cdot P' \rangle^{(-1)^{q+1}} \otimes \langle P \cdot \sfrak_{q-1}(L) \cdot \sfrak_1(L) \cdot P' \rangle^{(-1)^{q}} \simeq \Ocal_S.
\end{equation}
This identifies the leftmost bundle with the trivial line functor $\Ocal_S$. For the general case, one argues by applying the splitting principle in Proposition \ref{Prop:2ndsplitting} and the Whitney isomorphism above. This reduces us to the case of a product of Chern classes of line bundles with one or more indices strictly bigger than one. We can use the above isomorphism \eqref{eq:Segre-manipulation} for any of those Chern classes. It is formal to verify that the resulting trivialization does not depend on the order, by an argument similar to that after \eqref{eq:definition-chern-classes}. We notice that, in particular, the compatibility with the Whitney isomorphism is part of the construction of the rank triviality isomorphism.

The restriction isomorphism can be constructed as in \cite[\textsection V.4.9]{Elkikfib}. Set $T$ to be the image of $Y$ in $S$. It is both open and closed since $Y\to S$ is proper, flat, and of finite presentation. Since the construction of intersection bundles is local with respect to $S$, and by the convention as in Remark \ref{rmk:clarification-restriction}, we can suppose that either $S = T$ or that~$Y$ is empty.

We first suppose that $S=T$. We observe that, since $\sigma$ is a regular section whose zero locus is faithfully flat over $S$, then $Y$ automatically satisfies the condition $(C_{n-r})$ over $S$. Indeed, to verify that the morphism $Y \to S$ has constant pure relative dimension $n-r$, one can argue locally on $X$, in which case $Y \subseteq X$ is obtained by successively intersecting with relative Cartier divisors which are of pure relative dimension $n-1$, as in the remarks following Proposition \ref{prop:generalpropertiesdeligneproducts}.

We follow \cite[\textsection V.4.10]{Elkikfib}. The case $r=1$ relies on \eqref{item:deligne-pairing-restriction} in Proposition \ref{prop:generalpropertiesdeligneproducts}. For general~$r$, we proceed by induction on the rank. Introduce $\pi\colon\PBbb(E)\to X$, and the universal exact sequence \eqref{eq:tautological-exact-sequence}.

The section $\pi^{\ast}\sigma$ of $\pi^{\ast}E$ induces a regular section of $\Ocal(1)$ on $\PBbb(E)$. We denote by $D\subseteq\PBbb(E)$ its divisor. It is a relative Cartier divisor. The restriction of $\pi^{\ast}\sigma$ on $D$ defines a regular section of $Q|_{D}$. Its zero locus scheme $Z$ equals the restriction of the projective bundle to $Y$, that is $Z=\PBbb(E|_Y)$. Then there is a string of isomorphisms:
\begin{equation}\label{eq:definitionrestrictionisomorphism}
\begin{split}
\langle P \cdot \cfrak_r(E) \cdot P' \rangle_{X/S} & \simeq \langle \pi^\ast P \cdot \cfrak_r(\pi^\ast E) \cdot \pi^\ast P' \cdot \cfrak_{1}(\Ocal(1))^{r-1} \rangle_{\PBbb(E)/S} \\
& \simeq \langle \pi^\ast P \cdot \cfrak_{r-1}(Q)\cdot \cfrak_{1}(\Ocal(1)) \cdot \pi^\ast P' \cdot \cfrak_{1}(\Ocal(1))^{r-1} \rangle_{\PBbb(E)/S} \\
& \simeq \langle \pi^{\ast}P|_{D} \cdot \cfrak_{r-1}(Q|_D) \cdot \pi^\ast (P'|_D )\cdot \cfrak_{1}(\Ocal(1))^{r-1} \rangle_{D/S} \\
& \simeq \langle \pi^{\ast}P|_{Z} \cdot \pi^\ast (P'|_Z )\cdot \cfrak_{1}(\Ocal(1))^{r-1} \rangle_{Z/S} \\
& \simeq \langle P|_{Y} \cdot P'|_Y \rangle_{Y/S}.
\end{split}
\end{equation}
The first isomorphism in \eqref{eq:definitionrestrictionisomorphism} is the second case of the projection formula. The second one is the Whitney isomorphism associated with \eqref{eq:tautological-exact-sequence} combined with the rank triviality. The third one uses the case $r=1$, and the fourth isomorphism is constructed by induction. The final one performs the projection formula again.

The argument for the case $Y=\emptyset$ goes along the same lines. The case of rank one again reduces to Proposition \ref{prop:generalpropertiesdeligneproducts}. See Remark \ref{rmk:clarification-restriction}. For the higher rank case, the same manipulation as in \eqref{eq:definitionrestrictionisomorphism} applies, but here instead the scheme $Z$ will be empty, allowing us to apply induction again.

Finally, we address the birational invariance. We can suppose $P$ is a single monomial in the Segre classes of some vector bundles. If $p: \mathbf{P} \to X$ denotes the projective bundle over $X$ that appears in the definition of the Segre classes in $P$, then $\mathbf{P} \times_X \widetilde{X} \to \widetilde{X}$ is the corresponding projective bundle in the definition of the Segre classes of $\pi^\ast P$. The open subset $p^{-1}(U)$ satisfies the assumptions of Proposition \ref{prop:elvira} and the property easily follows from this.
\end{proof}

\begin{equation}\label{eq:tautological-exact-sequence}
0 \to Q \to \pi^\ast E \to \Ocal(1) \to 0.
\end{equation}

\begin{thebibliography}{99}
\bibitem{BDRR} N. A. Baas, B. I. Dundas, B. Richter \& J. Rognes - “Ring completion of rig categories” , J. reine angew. Math. 674 (2013), p. 43-80
\bibitem{freed1} J.-M. Bismut \& D. Freed - “The analysis of elliptic families. I. Metrics and connections on determinant bundles” , Comm. Math. Phys. 106 (1986) no. 1, p. 159-176
\bibitem{freed2} J.-M. Bismut \& D. Freed - “The analysis of elliptic families. II. Dirac operators, eta invariants, and the holonomy theorem” , Comm. Math. Phys. 107 (1986) no. 1, p. 103-163
\bibitem{BGS1} J.-M. Bismut, H. Gillet \& C. Soulé - “Analytic torsion and holomorphic determinant bundles. I. Bott-Chern forms and analytic torsion” , Comm. Math. Phys. 115 (1988) no. 1, p. 49-78
\bibitem{BGS2} J.-M. Bismut, H. Gillet \& C. Soulé - “Analytic torsion and holomorphic determinant bundles. II. Direct images and Bott-Chern forms” , Comm. Math. Phys. 115 (1988) no. 1, p. 79-126
\bibitem{BGS3} J.-M. Bismut, H. Gillet \& C. Soulé - “Analytic torsion and holomorphic determinant bundles. III. Quillen metrics on holomorphic determinants” , Comm. Math. Phys. 115 (1988) no. 2, p. 301-351
\bibitem{Borel-Serre} A. Borel \& J.-P. Serre - “Le théorème de Riemann-Roch” , Bull. Soc. math. France 86 (1958), p. 97-136
\bibitem{Boucksom-Eriksson} S. Boucksom \& D. Eriksson - “Spaces of norms, determinant of cohomology and Fekete points in non-Archimedean geometry” , Adv. Math. 378 (2021), article ID 107501, 124 pages
\bibitem{SebWalMarDiff} S. Boucksom, W. Gubler \& F. Martin - “Differentiability of relative volumes over an arbitrary non-Archimedean field” , Internat. Math. Res. Notices (2022) no. 8, p. 6214-6242
\bibitem{Deligne-determinant} P. Deligne - “Le déterminant de la cohomologie” , in Current trends in arithmetical algebraic geometry (Arcata, Calif., 1985) , Contemp. Math. , vol. 67 , American Mathematical Society, Providence, RI, 1987, p. 93-177
\bibitem{Ducrot} F. Ducrot - “Cube structures and intersection bundles” , J. Pure Appl. Algebra 195 (2005) no. 1, p. 33-73
\bibitem{Elkikfib} R. Elkik - “Fibrés d’intersections et intégrales de classes de Chern” , Ann. Sci. École Norm. Sup. (4) 22 (1989) no. 2, p. 195-226
\bibitem{Elkikmetr} R. Elkik - “Métriques sur les fibrés d’intersection” , Duke Math. J. 61 (1990) no. 1, p. 303-328
\bibitem{EriComptes} D. Eriksson - “Un isomorphisme de type Deligne-Riemann-Roch” , Comptes Rendus Mathématique 347 (2009) no. 19-20, p. 1115-1118
\bibitem{Eriksson-Freixas-Wentworth} D. Eriksson, G. Freixas i Montplet \& R. A. Wentworth - “Complex Chern–Simons bundles in the relative setting” , 2021
\bibitem{FrankeChow} J. Franke - “Chow categories” , Compositio Math. 76 (1990) no. 1-2, p. 101-162, Algebraic geometry (Berlin, 1988)
\bibitem{FrankeChern} J. Franke - “Chern functors” , in Arithmetic algebraic geometry (Texel, 1989) , Progress in Math. , vol. 89 , Birkhäuser Boston, Boston, MA, 1991, p. 75-152
\bibitem{Franke} J. Franke - “Riemann–Roch in functorial form” (1992), Unpublished
\bibitem{Fulton} W. Fulton - Intersection theory , Ergeb. Math. Grenzgeb. (3) , vol. 2 , Springer-Verlag, Berlin, 1998
\bibitem{FultonLang} W. Fulton \& S. Lang - Riemann-Roch algebra , Grundlehren Math. Wissen. , vol. 277 , Springer-Verlag, New York, 1985
\bibitem{Gabriel-Zisman} P. Gabriel \& M. Zisman - Calculus of fractions and homotopy theory , Ergeb. Math. Grenzgeb. (2) , vol. 35 , Springer-Verlag New York, Inc., New York, 1967
\bibitem{Munoz} E. Muñoz García - “Fibrés d’intersection” , Compositio Math. 124 (2000) no. 3, p. 219-252
\bibitem{GS-ARR} H. Gillet \& S. Soulé - “An arithmetic Riemann-Roch theorem” , Invent. Math. 110 (1992) no. 3, p. 473-543
\bibitem{JG} P. Goerss \& J. Jardine - Simplicial homotopy theory , Progress in Math. , vol. 174 , Birkhäuser Verlag, Basel, 1999
\bibitem{Grayson} D. Grayson - “Higher algebraic K -theory. II (after Daniel Quillen)” , in Algebraic K -theory (Proc. Conf., Northwestern Univ., Evanston, Ill., 1976) , Lect. Notes in Math. , vol. 551 , Springer, 1976, p. 217-240
\bibitem{EGAII} A. Grothendieck - “Éléments de géométrie algébrique. II. Étude globale élémentaire de quelques classes de morphismes” , Publ. Math. Inst. Hautes Études Sci. 8 (1961), p. 5-222
\bibitem{EGAIV1} A. Grothendieck - “Éléments de géométrie algébrique. IV. Étude locale des schémas et des morphismes de schémas. I” , Publ. Math. Inst. Hautes Études Sci. 20 (1964), p. 5-251
\bibitem{EGAIV2} A. Grothendieck - “Éléments de géométrie algébrique. IV. Étude locale des schémas et des morphismes de schémas. II” , Publ. Math. Inst. Hautes Études Sci. 24 (1965), p. 5-231
\bibitem{EGAIV3} A. Grothendieck - “Éléments de géométrie algébrique. IV. Étude locale des schémas et des morphismes de schémas. III” , Publ. Math. Inst. Hautes Études Sci. 28 (1966), p. 5-255
\bibitem{EGAIV4} A. Grothendieck - “Éléments de géométrie algébrique. IV. Étude locale des schémas et des morphismes de schémas IV” , Publ. Math. Inst. Hautes Études Sci. 32 (1967), p. 5-361
\bibitem{SGA6} - Théorie des intersections et théorème de Riemann-Roch (SGA 6) , Lect. Notes in Math. , vol. 225 , Springer-Verlag, Berlin-New York, 1971, Séminaire de Géométrie Algébrique du Bois-Marie 1966–1967 (SGA 6), Dirigé par P. Berthelot, A. Grothendieck et L. Illusie. Avec la collaboration de D. Ferrand, J. P. Jouanolou, O. Jussila, S. Kleiman, M. Raynaud et J. P. Serre
\bibitem{Johnson-Yau} N. Johnson \& D. Yau - “Bimonoidal categories, E n -monoidal categories, and algebraic K -theory” , 2021
\bibitem{Knudsen} F. Knudsen - “Determinant functors on exact categories and their extensions to categories of bounded complexes” , Michigan Math. J. 50 (2002) no. 2, p. 407-444, Erratum: Ibid. 50 (2002), no. 2, p. 407–444
\bibitem{KnudsenMumford} F. Knudsen \& D. Mumford - “The projectivity of the moduli space of stable curves. I. Preliminaries on “det” and “Div”” , Math. Scand. 39 (1976) no. 1, p. 19-55
\bibitem{Laplaza} M. L. Laplaza - “Coherence for distributivity” , in Coherence in categories , Lect. Notes in Math. , vol. 281 , Springer-Verlag, Berlin-Heidelberg-New York, 1972, p. 29-65
\bibitem{LipNee} J. Lipman \& A. Neeman - “Quasi-perfect scheme-maps and boundedness of the twisted inverse image functor” , Illinois J. Math. 51 (2007) no. 1, p. 209-236
\bibitem{MacLane} S. Mac Lane - “Natural associativity and commutativity” , Rice Univ. Stud. 49 (1963) no. 4, p. 28-46
\bibitem{MacLane-cat} S. Mac Lane - Categories for the working mathematician , Graduate Texts in Math. , vol. 5 , Springer-Verlag, New York-Berlin, 1971
\bibitem{Mumford-red} D. Mumford - The red book of varieties and schemes , Lect. Notes in Math. , vol. 1358 , Springer-Verlag, Berlin, 1999
\bibitem{Muro:determinant} F. Muro, A. Tonks \& M. Witte - “On determinant functors and K -theory” , Publ. Mat. 59 (2015) no. 1, p. 137-233
\bibitem{Nakai} Y. Nakai - “Some fundamental lemmas on projective schemes” , Trans. Amer. Math. Soc. 109 (1963), p. 296-302
\bibitem{Nakayamaintsheaves} N. Nakayama - “Intersection sheaves over normal schemes” , J. Math. Soc. Japan 62 (2010) no. 2, p. 487-595
\bibitem{Patel} D. Patel - “de Rham $\mathcal{E}$ -factors” , Invent. Math. 190 (2012) no. 2, p. 299-355
\bibitem{Quillen:K-theory-I} D. Quillen - “Higher algebraic K -theory. I” , Algebraic K -theory, I: Higher K -theories (Proc. Conf., Battelle Memorial Inst., Seattle, Wash., 1972) (Lect. Notes in Math.) 341 , Sprinegr (1973), p. 85-147
\bibitem{Richter} B. Richter - From categories to homotopy theory , Cambridge Studies in Advanced Math. , vol. 188 , Cambridge University Press, Cambridge, 2020
\bibitem{Rossler-ARR} D. Rössler - “A local refinement of the Adams-Riemann-Roch theorem in degree one” , in Arithmetic L-functions and differential geometric methods , Progress in Math. , vol. 338 , Birkhäuser/Springer, Cham, 2021, p. 213-246
\bibitem{Saavedra} N. Saavedra Rivano - Catégories tannakiennes , Lect. Notes in Math. , vol. 265 , Springer-Verlag, Berlin-New York, 1972
\bibitem{Sinh-these} H. X. Sính - Gr -catégories , Ph. D. Thesis, Institut pédagogique no. 2 de Hanoi, 1975,
\bibitem{Suslin-Voevodsky} A. Suslin \& V. Voevodsky - “Relative cycles and Chow sheaves” , in Cycles, transfers, and motivic homology theories , Annals of Math. Studies , vol. 143 , Princeton Univ. Press, Princeton, NJ, 2000, p. 10-86
\bibitem{stacks-project} The Stacks Project Authors - “ Stacks Project ” , , 2023
\bibitem{ThomasonPhony} R. Thomason - “Beware the phony multiplication on Quillen’s $\mathcal{A}$ - 1 $\mathcal{A}$ ” , Proc. Amer. Math. Soc. 80 (1980) no. 4, p. 569-573
\bibitem{ThomasonSpectral} R. Thomason - “First quadrant spectral sequences in algebraic K -theory via homotopy colimits” , Comm. Algebra 10 (1982) no. 15, p. 1589-1668
\bibitem{ThomasonHomCol} R. W. Thomason - “Homotopy colimits in the category of small categories” , Math. Proc. Cambridge Philos. Soc. 85 (1979) no. 1, p. 91-109
\bibitem{ThomasonTrobaugh} R. W. Thomason \& T. Trobaugh - “Higher algebraic K -theory of schemes and of derived categories” , in The Grothendieck Festschrift, Vol. III , Progress in Math. , vol. 88 , Birkhäuser Boston, Boston, MA, 1990, p. 247-435
\bibitem{Waldhausen} F. Waldhausen - “Algebraic K -theory of spaces” , in Algebraic and geometric topology (New Brunswick, NJ, 1983) , Lect. Notes in Math. , vol. 1126 , Springer, 1985, p. 318-419
\end{thebibliography}
\end{document}
